% ============================================================
% Tabelle 5.3-1: Verarbeitungsschritte des Evidence Ledgers — v01 (09.09.)
% Label: t:ledger-schritte
% Stil: chap6.1-Konvention.
% HERKUNFT: Reihenfolge = tatsächliche Graphverdrahtung
% LedgerBuilderWorkflow.cs (Builder „LedgerBuilderUnitCoverage":
% Segmentation → Extraction → UnitCoverageGate → UnusedTriage →
% UnusedCompare → CompareRepair → Canonicalization → [Check ⇄ Repair]
% → FacetValidation); Art je Schritt = F-1/F-2 des Umsetzungsplans
% (FacetValidator = LLM-Call, det. Antwortprüfung + Statusableitung
% grounded→approved; UnusedUnitCompareRepair = det. Referenzpflicht-
% Abstufung); Queue-Filter = LedgerAdjudicationAdapter
% (DecidableCompareVerdicts); Leer-Queue-Skip = R-50
% (AdjudicationStageExecutor). Funktional gebündelt, kein
% Executor-Inventar.
% NACHZUG 09.09.: letzte Zeile = A10-Verfeinerung
% (PipelineAdjudicationRefine, geteilte Naht AdjudicationRefine.cs) —
% am 09.09. nach Abschluss der Messkampagne eingebunden
% (Standabgrenzung in 5.6; Kap.-7-Läufe entstanden DAVOR).
% ============================================================

\begin{table}[htbp]
  \centering
  \small
  \begin{tabular}{>{\raggedright\arraybackslash}p{0.235\textwidth}>{\raggedright\arraybackslash}p{0.135\textwidth}>{\raggedright\arraybackslash}p{0.51\textwidth}}
    \hline
    \textbf{Schritt (Eingabe $\rightarrow$ Ergebnis)} & \textbf{Art} & \textbf{Fortsetzung / menschliche Klärung} \\
    \hline
    Unit-Segmentierung (Transkript $\rightarrow$ adressierbare Units) & deterministisch & immer weiter \\
    Kandidaten-Extraktion (Units $\rightarrow$ quellgebundene Claim-Kandidaten) & modellgestützt & weiter \\
    Unit-Bilanz (verwendete / unverwendete Units) & deterministisch & weiter; Grundlage der Hinweisprüfung \\
    Hinweisprüfung ungenutzter Units (Triage und Vergleich gegen die Kandidaten) & modellgestützt; anschließend deterministische Belegpflicht & Prüfergebnisse je Unit (fehlender Claim · menschlich zu klären · als Beleg anzuhängen · bereits mittelbar abgedeckt); ein Deckungsurteil ohne benannten Kandidaten wird deterministisch zu „menschlich zu klären`` abgestuft \\
    Kanonisierung (Kandidaten $\rightarrow$ kanonischer Claim-Bestand, Bündelung möglich) & modellgestützt & deterministischer Übergangs-Check prüft die \emph{referenzielle} Deckung (jeder Kandidat referenziert?); bei Verstoß Reparaturschleife, sonst weiter \\
    Facetten-Validierung (Facetten und Dispositionen je Claim) & Modell-Urteil; deterministische Antwortprüfung und Statusableitung & als gestützt beurteilte Claims werden freigegeben; alle übrigen werden zur menschlichen Nachprüfung markiert \\
    Vorlagen-Aufbau (Adjudikations-Liste) & deterministisch & markierte Claims und Hinweise werden vorgelegt; „bereits mittelbar abgedeckt``-Hinweise werden nicht vorgelegt \\
    Adjudikation und Anwendung & Mensch; deterministische Anwendung & Entscheidung je Position; Ergebnis ist die verwendbare Evidenzsicht --- eine leere Vorlagenliste überspringt den Entscheidungspunkt \\
    Facetten-Vervollständigung neuer Claims & Modell-Urteil (nur bei ausstehenden Facetten) & vervollständigt Facetten und Dispositionen per Adjudikation neu aufgenommener Claims vor der fachlichen Ableitung; ohne ausstehende Facetten kein Modellaufruf (Durchreich-Schritt) \\
    \hline
  \end{tabular}
  \caption[Verarbeitungsschritte des Evidence Ledgers]{Die tatsächlich realisierten Verarbeitungsschritte des
    Evidence Ledgers in Ausführungsreihenfolge (funktional gebündelt).
    „Modellgestützt`` bezeichnet ein semantisches Modellurteil;
    deterministische Schritte prüfen Form, Referenzen oder leiten
    Status ab. Die Hinweisprüfung arbeitet auf dem
    \emph{Kandidaten}-Bestand vor der Kanonisierung.}
  \label{t:ledger-schritte}
\end{table}
